\section{Interoperability and Interpretability}
\label{sec:interop-interpret}
Two requirements recur across every plan-of-work goal but have no single home in a fusion
architecture on their own: the system must accept whatever subset of a patient's data streams is
actually available (interoperability), and every number it emits must be traceable to why it was
produced (interpretability). We treat both as first-class, machine-checked properties rather than
narrative claims.

\subsection{Interoperability}
\label{subsec:interop}
\paragraph{Device and record boundaries.} Real-time ingestion from Galea, EMOTIV, Muse, and
Empatica devices is defined by explicit Lab Streaming Layer (LSL) stream contracts spanning three
logical modalities --- EEG, wearable physiology, and reference glucose --- captured via
LabRecorder and replayable offline for testing. Outbound, every real-time frame carries not a bare
value but its confidence and an explicit \texttt{evidence\_status}: the system abstains rather than
emitting a number when the input is insufficient. Clinical-record interoperability is handled the
same way: a glucose forecast can be exported as a LOINC-coded FHIR \texttt{Observation} resource
carrying its own prediction interval and research-only status, so a forecast enters an electronic
health record as a structured, interpretable object rather than a loose number.

\paragraph{Missing modalities, formalized.} Section~\ref{subsec:irm}'s IRM mixer makes the same
principle provable at the model level rather than only at the interface level: because each extra
modality enters as an independent gated residual (Equation~\ref{eq:irm-mix}), dropping one changes
the fused representation by exactly that modality's own gated contribution
(Equation~\ref{eq:irm-bound}) --- a closed-form bound that Kumar et al.'s convex mix and the
sequential updates of comparable modular architectures~\cite{multimodn} do not offer. The
consistency regularizer of Section~\ref{subsec:irm} is a challenger option, off by default; this paper does not treat a cross-cohort glycemic board as a parent result.

\paragraph{Invariants enforced everywhere.} Three invariants hold identically across the real-time
stream, the FHIR export, and the REST prediction API: a risk band and, where applicable, a grounded
reason accompany every probability; a missing-data state is transmitted as an explicit abstention,
never a fabricated value or a silent zero; and \texttt{validated\_for\_clinical\_use=False} with a
research-stage disclaimer is attached at every boundary, so no downstream consumer can mistake a
research estimate for a cleared clinical result. Serving-path latency stays real-time-safe
throughout (direct path 0.15\,ms, agent-orchestrated path 2.3\,ms, real-time frame path
${\sim}$0\,ms; glucose point models predict in under 1.5\,ms).

\subsection{Interpretability}
\label{subsec:interpret}
Every prediction is explained through three complementary, machine-checked layers, none of which
invents a number (\texttt{tests/test\_no\_hallucinated\_numbers.py} enforces this directly).

\paragraph{Signed feature attributions.} Every response carries per-feature signed contributions
from a shared linear-attribution surface, e.g., for a glucose-instability prediction: time-above-range
and CGM variability raise the estimate, HbA1c and fasting glucose raise it further --- the
consumer sees which inputs pushed the estimate up or down, and by how much, rather than a black
box.

\paragraph{Model-intrinsic importances.} The gradient-boosted and tree models in the glucose model
ladder expose native feature importances; the CGM lag and slope features they rank highest are the
same causal-history signals the deep model consumes, cross-checking that the linear-attribution
explanation is faithful to the data rather than post-hoc storytelling.

\paragraph{Grounded natural-language explanation.} A narration layer restates only values already
frozen in the prediction body (\texttt{explanation.predicts == False} by construction) and is
confined to a single gated leaf in the serving graph; every numeric token in its text is asserted
to already appear in the prediction body it explains.

\paragraph{IRM-specific diagnostics.} For the attention-fusion pathway (Section~\ref{subsec:irm})
we additionally report, per prediction and per fold: the unique-information score $\Delta_k$
(how much task loss modality $k$'s presence actually removes) and its empirical distance
correlation~\cite{szekely2007dcor} with the primary pathway's representation --- together
distinguishing a modality that is redundant with what the model already knows from one that is
contributing genuinely independent signal. Both are diagnostic only and are never fed back into
training, preserving this project's frozen-specification discipline.
